% Measured source: artifacts/advancement/analysis/analysis.json
% SHA-256: 3727de01958ef96deb0c15ad22823506eb9097e975f046334d84976a10e37762
\subsection{Exact evaluation against released fast-VUS on two machines}
A new complete run compares the unchanged fast-VUS module at full thresholds
against Sparse VUS and the three dense implementations. Each machine evaluates
all 2,100 original model/width cases on 350 series from 23 families. Each call
includes validation, preprocessing, and integration; three repetitions use
randomized method order, one compute thread, and excluded JIT warmup.
The M1 Pro run uses Python 3.12.4, NumPy 1.26.4, and Numba 0.59.1.
The Linux hosted runner identifies an AMD EPYC 9V74 and uses Python 3.12.15,
NumPy 2.5.4, and Numba 0.68.0. Ordinary OS and hosted-hardware background load
is uncontrolled. The Linux run is a separate execution, not an external
laboratory reproduction.

Table~\ref{tab:fast-vus} reports within-machine ratios. Fast-VUS returns only
PR; Sparse VUS computes both ROC and PR. The per-case fastest-tested comparator
is the minimum median runtime among fast-VUS and all three dense evaluators,
excluding Sparse VUS itself. This is an oracle comparison, not a deployable
selector. The maximum PR difference from fast-VUS is $2.78\cdot10^{-15}$;
the largest discrepancy across all exact methods is $3.70\cdot10^{-13}$.

\begin{table}[t]
\centering
\begin{tabular}{llrrrr}
\toprule
Machine & Comparator & Median & \shortstack{Geometric mean\\(95\% interval)} & Faster & \shortstack{Slower\\or equal}\\
\midrule
M1 Pro & fast-VUS & 8.38 & 8.69 [8.32, 9.05] & 2100 & 0 \\
M1 Pro & Fastest tested & 2.04 & 1.98 [1.89, 2.07] & 1891 & 209 \\
EPYC 9V74 & fast-VUS & 4.55 & 4.50 [4.34, 4.68] & 2100 & 0 \\
EPYC 9V74 & Fastest tested & 1.35 & 1.40 [1.34, 1.46] & 1526 & 574 \\
\bottomrule
\end{tabular}
\caption{Comparator runtime divided by Sparse VUS runtime; values above one
favor Sparse VUS. Every row covers 2,100 cases. Intervals resample whole series
and describe geometric means on the fixed corpus. Both methods use all
thresholds; fast-VUS computes PR only.}
\label{tab:fast-vus}
\end{table}

Fast-VUS ratios range from 1.75 to 46.45 on M1 Pro and 1.52 to 17.97 on Linux.
Against the fastest-tested comparator, ranges are 0.40--7.75 and 0.36--8.40.
Whole-family bootstrap intervals for the latter geometric means are
[1.53,2.42] and [1.11,1.68], respectively. The corpus and machines support a
bounded measured advantage, with substantial dense-baseline slowdown cases;
they do not settle the unavailable Sylligardos comparison.
The sums of isolated per-case medians are 14.78 versus 104.52 seconds for
Sparse VUS and fast-VUS on M1 Pro, and 14.39 versus 55.14 seconds on Linux.
These sums are not measured batch or detector-pipeline wall times.

\subsection{Error bounds and the cost of width refinement}
The width-refinement experiment retains the metadata-selected 42 series
covering all 23 families and all three released score vectors. Maximum widths
30, 100, and 300 give 378 model/width cases; the width-300 setting is an
explicit stress condition. Both $\varepsilon=0.01$ and $0.001$ bound the
half-widths of both returned metric intervals. The M1 Pro timing protocol
again includes whole-call preprocessing and uses three randomized-order
repetitions. Table~\ref{tab:width-refinement} compares against full-width
Sparse VUS on these same inputs.

\begin{table}[t]
\centering
\begin{tabular}{rrrrrr}
\toprule
$\varepsilon$ & $L$ & \shortstack{Median\\ratio} & Faster / 126 & \shortstack{Distinct\\widths (\%)} & \shortstack{Curve calls\\(\%)}\\
\midrule
0.01 & 30 & 0.835 & 27 & 16.1 & 16.1 \\
0.01 & 100 & 0.905 & 46 & 7.9 & 7.9 \\
0.01 & 300 & 1.313 & 100 & 4.7 & 7.3 \\
0.001 & 30 & 0.615 & 4 & 54.8 & 80.6 \\
0.001 & 100 & 0.426 & 11 & 39.6 & 74.3 \\
0.001 & 300 & 0.391 & 25 & 25.4 & 53.8 \\
\bottomrule
\end{tabular}
\caption{Full-width exact runtime divided by adaptive runtime. Width and
curve-call columns are medians relative to $L+1$; cache eviction can make
curve calls exceed the number of distinct widths. Values below one indicate
that adaptive evaluation is slower.}
\label{tab:width-refinement}
\end{table}

All 378 cases at both tolerances contain the full-width ROC and PR values
within the separately declared $10^{-10}$ numerical audit allowance.
Maximum midpoint errors are $0.003260$ and $0.000713$ for tolerances 0.01
and 0.001. However, fewer width evaluations do not generally imply faster
execution. Across all three maximum widths, median runtime ratios are
0.966 and 0.467, with adaptive evaluation faster in 173 and 40 of 378 cases,
respectively. At tolerance 0.01 and $L=300$, the geometric mean ratio is
1.370 (series-bootstrap interval [1.207,1.557]); at the tighter tolerance,
all three width groups have geometric mean ratios below one.
The present implementation therefore provides useful enclosures, but the
standalone fixed-tolerance strategy is often slower than exact evaluation.

As a numerical decision replay, nonoverlapping PR intervals identify unique
winners in 21, 17, and 15 of the 42 series at tolerance 0.01 for widths
30, 100, and 300; the corresponding counts at tolerance 0.001 are 33, 32,
and 34. All asserted winners agree with exact PR using the declared
$2\cdot10^{-10}$ separation margin. These replay counts do not include
additional decision-workload timing or establish a model-selection speedup.

Separately authored validation directly constructs full-threshold curves
for every both-class binary label vector of lengths two through seven,
using four seeded score patterns per vector. It checks 13,776 widths and
width enclosures and 960 discrete averages. The largest curve discrepancy
is $6.67\cdot10^{-16}$ and the largest enclosure violation is
$3.34\cdot10^{-16}$, within the declared $2\cdot10^{-12}$ allowance.
A further 760 linear-programming comparisons of p-box extrema have maximum
discrepancy $3.34\cdot10^{-16}$. Exhaustiveness concerns labels only, not all
rankings or tie patterns; these checks audit the implementation and complement
the real-arithmetic proofs.

